\section{Introduction}

Financial QA requires evidence selection, arithmetic, and careful handling of
entities, periods, signs, and units. Iterative language-model agents appear
well suited to these demands: they can search, inspect evidence, and revise an
answer. Each extra step, however, consumes latency and tokens and creates
another opportunity for retrieval, arithmetic, or formatting error. The
deployment question is therefore not simply whether an agent is capable, but
whether it contributes answers that a cheaper workflow misses.

We study this necessary precondition for selective agency: \emph{expert
complementarity}. A router can improve an accuracy--cost frontier only when its
expensive branch supplies enough unique correct answers. This requirement is
logically prior to learning a difficulty model or tuning a threshold. It also
differs from asking which branch has higher average accuracy: two systems can
tie while making nearly identical errors, leaving little routing headroom.

Our main evidence is a prospectively frozen capability ladder on 150 fresh
development items from FinQA, TAT-QA, and ConvFinQA. We compare one-call Static
and bounded retrieval-first ReAct under a shared strict action contract,
retrieval interface, scorer, and resource ceilings using GPT-4o-mini, GPT-5.6
Luna, and GPT-5.6 Terra. Static significantly outperforms ReAct for the first
two tiers, whereas Terra reaches parity; the positive cross-tier interaction
establishes a capability boundary. Nevertheless, ReAct supplies only two, one,
and two unique exact successes, so oracle routing headroom remains at most 1.3
points in every tier.

An earlier five-dataset selector study and a separate harmonized comparison
motivate and support this analysis. The original selector failed its
quality--cost rule but used unequal answer and context contracts. The matched
study corrected those asymmetries and again found little ReAct
complementarity. We preserve both as supporting evidence rather than treating
the flawed pilot as the centerpiece.

Our contribution is thus a feasibility audit for routing: model capability can
change the Static--ReAct ordering, but it does not automatically create unique
expert value. Evaluation contracts must first be harmonized, then
complementarity measured, before selector optimization. This is a scoped result
for bounded retrieval-first ReAct over controlled benchmark evidence, not a
universal claim against agents.
